\bookchapter{Chapter 11: Photography and Image Management}

\chapterintro{You never adjust exposure}

\booksection{Have You Made These Mistakes Too?}

The flash is always on

\begin{itemize}
  \item \textbf{: The effect of a phone's LED flash is really ugly — the light is harsh and cold, making people look like ghosts. Unless it's truly pitch black, turn it off. In low light, Night Mode works far better than the flash.}Zooming by pinching with two fingers
  \item \textbf{: Pinch-to-zoom on the screen is digital zoom, which essentially crops and enlarges the image, noticeably reducing quality. If you want to shoot something far away, use the zoom buttons on the viewfinder (0.5x/1x/2x/3x) — these are optical zoom and don't lose quality.}2. A Few Small Tips to Immediately Boost Your Hit Rate
  \item \textbf{These don't require any theoretical knowledge — use them once and you'll see the effect:}Rule of Thirds + Grid Lines
  \item \textbf{: Turn on "Grid Lines" (also called reference lines) in camera settings. When shooting, place what you want to highlight at the intersections of the four lines, and the composition will immediately "click". This is the simplest and most effective rule in photography — no mysticism, it just matches the visual habits of the human eye.}: Pinch-zooming on the screen is digital zoom — essentially cropping and enlarging the image, which noticeably reduces quality. To zoom in on distant subjects, use the zoom buttons in the viewfinder (0.5x/1x/2x/3x). The 0.5x (ultrawide) and 1x (main camera) are native optical lenses with no quality loss; 2x/3x are typically "lossless quality zoom" — achieved by cropping from high-resolution sensors, with near-optical quality but not a separate lens. Beyond 3x is usually hybrid or digital zoom, with noticeable quality degradation.
\end{itemize}

\booksection{A Few Small Tips to Immediately Boost Your Hit Rate}

Don't always shoot standing up

\begin{itemize}
  \item \textbf{: Most people always shoot at their own eye level, so their photos all have the same perspective. Try crouching down for a low angle, or finding a high spot to shoot downward — the same place can produce a completely different feel. When shooting kids and pets especially, you need to crouch down — only eye-level shots produce a natural feel.}Light is the top priority
  \item \textbf{: Shooting portraits under the bright midday sun produces hard, ugly shadows on the face. The best light is during the "golden hours" of early morning and late evening — the sunlight is angled and warm-toned, making everything look good. Overcast-day light isn't flashy, but it's even and soft, making it actually easier to shoot portraits than under bright sun.}Night Mode essentials
  \item \textbf{: After turning on Night Mode, the phone automatically takes multiple shots and composites them. The most critical thing at this point is — your hand absolutely cannot shake. Lean against a wall if you can, rest on a table if you can, or hold your breath if you must. Night Mode requires you to hold steady for 2 to 4 seconds — any shake and the image blurs.}3. Portrait Mode and Pro Mode
  \item \textbf{Portrait Mode}: It produces the effect of a blurred background with a sharp subject (shallow depth of field). Because phones have small lens sizes, their real bokeh capability is limited, so this relies on AI algorithms — simply put, the AI cuts you out and blurs the background. The best results come at a distance of 1-2 meters; too far or too close and the algorithm easily makes mistakes. Also, when shooting portraits, make sure the background doesn't have overly complex lines, or the AI cutout may "fail" around the hair edges.
  \item \textbf{Pro Mode}: If you want full control over the image effect, try Pro Mode. A few key parameters to know: the lower the ISO, the cleaner the image quality, but the darker the image; the slower the shutter, the more light comes in, but hand shake causes blur; white balance controls color temperature — set poorly, photos will look blue or yellow. Most people don't need to use Pro Mode year-round, but understanding these parameters at least lets you know what the camera is doing.
\end{itemize}

\booksection{Portrait Mode and Pro Mode}

\begin{itemize}
  \item \textbf{Many people misunderstand photo editing, thinking that editing makes it "unreal". In fact, photos taken by a camera are already different from what the naked eye sees — appropriate adjustment is to restore the scene you saw at the time, not to fake it.}Snapseed
  \item \textbf{: A free editing tool from Google, professional features with no ads. My daily editing flow is: first use "Tune Image" to adjust brightness and contrast use "Details" to sharpen use "White Balance" to correct color casts finally "Crop" to the desired composition. Snapseed's "Selective" feature is very useful — you can adjust only a specific area of the image without affecting the whole thing.}Xingtu / Meitu
\end{itemize}

\booksection{Photo Editing: Don't Overdo It — Just Right Is Enough}

Lightroom Mobile

\begin{itemize}
  \item \textbf{: If you're particular about color, this is the most professional. It has advanced features like curve adjustment and HSL per-color adjustment. But for most people, it's overkill.}Editing Workflow Suggestion
  \item \textbf{: My own habit is to shoot the original photo first, then edit slowly on a computer screen or iPad at home — the phone screen is too small and the color isn't quite accurate. If you must edit on the phone, remember to set the screen brightness to medium-high, and don't edit in power-saving mode.}5. Too Many Photos? Organization Strategies
  \item \textbf{Thousands of photos piled up in the phone album, taking forever to find a specific one — everyone has had this experience. Here are a few practical organization methods:}Make good use of search
  \item \textbf{: Now both iPhone and most Android phone albums support smart search. You can search for keywords like "cat", "seaside", "ID card", "flowers", and the system will use AI to identify image content and find them for you. You can also search by date and location, like "last May" or "Beijing". Not many people know about this feature, but it's really useful.}Regularly clean up junk
\end{itemize}

\booksection{Too Many Photos? Organization Strategies}

Create albums — don't rely solely on default sorting

\begin{itemize}
  \item \textbf{: The system sorting by time is fine, but for cross-time content like "child's growth photos", "travel photos", and "important documents", you should manually create albums to categorize them. Apple's "Shared Albums" can even let family members add photos together — one person creates it, and the whole family can view it.}Backup is the bottom line
  \item \textbf{: Photos are the most irreplaceable data — files can be recreated, chat histories can be lost and that's that, but a child's photos or an elderly person's videos, once lost, are truly gone. You must enable cloud backup. iCloud Photos is the most worry-free; on Android, each brand has its own cloud service, or you can use Baidu Netdisk / Aliyun Drive for auto backup. The principle is: have at least one cloud backup — don't put all your eggs in one basket.}6. Screenshots and Screen Recording
  \item \textbf{Screenshots}: Press the power button + volume up button simultaneously — a universal method for all phones. Many Android phones also support three-finger swipe down screenshots and knuckle double-tap screenshots. Long screenshots (scrolling screenshots) are well-established on Android — after taking one shot, tap the "Scroll" button, and the phone automatically scrolls down and stitches together a long image, suitable for capturing chat histories and web pages.
  \item \textbf{Screen Recording}: Add "Screen Recording" from the Control Center, tap once to start recording, tap again to stop. Before recording, long-press the record button to choose whether to record microphone sound — turn on the mic if you're making a tutorial with narration, turn it off if you're just recording operations. The recorded video is automatically saved to the album.
\end{itemize}

\booksection{Screenshots and Screen Recording}

\begin{itemize}
  \item \textbf{Turn on "Grid Lines" in camera settings, find a spot with natural light, and take three photos (remember to wipe the lens and adjust exposure).}Select the best one, and use Snapseed for basic adjustments (brightness, contrast, cropping).
  \item \textbf{Create a dedicated album for this photo, named "My Works".}Check whether cloud backup is enabled, to ensure the photo has a copy in the cloud.
\end{itemize}
